\documentclass[12pt]{article}

\input{../../../_assets/preambulo}


\begin{document}

    % 1. Foto de fondo
    % 2. Título
    % 3. Encabezado Izquierdo
    % 4. Color de fondo
    % 5. Coord x del titulo
    % 6. Coord y del titulo
    % 7. Fecha

    
	\input{../../../_assets/portada}
	\portadaExamenFotoDif{../../sapienza.jpg}{Analisi Funzionale\\Examen I}{Analisi Funzionale. Examen I}{MidnightBlue}{-8}{28}{2025-2026}{Joaquín Avilés de la Fuente}

    \begin{description}
        \item[Asignatura] Analisi Funzionale.
        \item[Curso Académico] 2025-2026.
        \item[Grado] Grado en Matemáticas.
        \item[Grupo] Grupo Único.
        \item[Profesor] Filomena Pacella.
        \item[Descripción] Examen final escrito de Analisi Funzionale.
        \item[Fecha] 4 de junio de 2026.
        \item[Duración] 150 min.
    
    \end{description}
    \newpage

    % ------------------------------------
    
    \begin{ejercicio}[4 punti]
        Dimostrare che se $X$ è uno spazio di Banach riflessivo allora:
		\[
		\forall F \in X^* \quad \exists \, \widetilde{x} \in X : \|\widetilde{x}\| = 1 \quad \text{e} \quad \|F\| = |F(\widetilde{x})|
		\]


    \end{ejercicio}

	\begin{ejercicio}[5 punti]
		Siano $X$ e $Y$ spazi normati e $T : X \to Y$ une applicazione lineare e chiusa. Dimostrare:
		\begin{enumerate}
			\item[(i)] se $T$ è biettiva allora 
			$$T^{-1} : Y \to X$$ 
			è lineare e chiusa.
			\item[(ii)] se $K \subseteq X$ è compatto allora $T(K)$ è chiuso in $Y$.
		\end{enumerate}
	\end{ejercicio}

	\begin{ejercicio}[4 punti]
		Siano $X$ e $Y$ spazi di Banach e $T \in BL(X,Y)$. Dimostrare che se $T$ è compatto allora:
		\[
		x_n \rightharpoonup x_0 \text{ in } X \implies T(x_n) \to T(x_0)
		\]
	\end{ejercicio}

    \begin{ejercicio}[4 punti]
		Sia $H$ uno spazio di Hilbert di dimensione infinita e 
		$$T : H \to H$$
		un operatore lineare e compatto. Dimostrare che $\not\exists \, \alpha > 0$ tale che:
		\[
		(T(x), x) \ge \alpha \|x\|^2, \qquad \forall x \in H
		\]
	\end{ejercicio}

	\begin{ejercicio}[8 punti]
		Sia $X$ uno spazio di Banach e $T \in BL(X)$. Definiamo $\lambda \in \mathbb{R}$ autovalore approssimato di $T$ se: \\
		$$\exists \{x_n\} \subset X : \|x_n\| = 1, \ \forall n \in \mathbb{N}, \ \text{ e } \ T(x_n) - \lambda x_n \to 0, \ n \to \infty$$
		Dimostrare che:
		\begin{enumerate}
			\item[(i)] ogni autovalore di $T$ è anche autovalore approssimato di $T$.
			\item[(ii)] ogni autovalore approssimato di $T$ appartiene allo spettro di $T$.
		\end{enumerate}		
	\end{ejercicio}

	\begin{ejercicio}[8 punti]
		Sia $T : C^0([1, e]) \to C^0([1, e])$ definito da:
		\[
		T(u)(t) = (\log t) \, u(t), \qquad \forall t \in [1, e], \ \forall u \in C^0([1, e])
		\]
		Determinare lo spettro di $T$, gli eventuali autovalori e i relativi autospazi. L'operatore $T$ è compatto?
	\end{ejercicio}
\end{document}